            In the case of Zynq the interface uses AXI bus GP0.            
            The size of the output data type must not exceed 256 bits as the CPU
            interface software limits burst transfers to 8 32-bit words.
            
            Boards having a serial interface to the CPU use that interface.
            The size of the output data type is not limited.

            The queue must be read somewhere in the program (queues without
            sinks result in a fatal error on completion of immediate execution).
            When the module is created the argument queue must have an established
            read clock domain or a fatal error is generated.

            If the output queue is written anywhere else in the program (i.e. other
            than in this module) then it must be written in the bus clock domain
            (c100) since that is the clock domain in which it is written in this
            module.

            If the argument to {\bf freecount} is true the number of free entries in
            the queue buffer can be interrogated by the CPU. The output queue may
            have any of the allowed buffer sizes. If the queue is asynchronous, i.e.
            the write domain is not the bus clock domain (c100), the count will not
            be exact but will be a guaranteed minimum number of writable spaces.

            If the argument to {\bf freecount} is false the data returned from
            reading the number of free spaces in the buffer is 0 or 1 instead of a
            count. A return of 1 means the queue buffer is not full, i.e. can be
            written. This option is provided to simplify the logic where a simple
            full/non-full status is required.

            If {\bf freecount} has no argument then there will no interface to
            directly read the space remaining in the queue buffer, however an
            interrupt may be used instead.            

            Arguments to parameter {\bf ilevel} control provision of a level
            interrupt on queue buffer status as follows -

            If there is no argument then no interrupt is provided.

            If the argument is type "log" with value false then the interrupt will
            be asserted when the queue buffer is not full.

            If the argument is type "log" with value true then the interrupt will be
            asserted when the number of free spaces in the buffer equals or exceeds
            a level which has been written to the interface from the CPU. The
            initial value of this level is 1. A level value of 0 should not be
            written to the interface as that will result in a continuous level
            interrupt regardless of queue state.

            If the argument is type "uint" or "int" then the interrupt will be
            asserted when the number of free spaces in the buffer equals or exceeds
            the fixed value specified by the integer argument. Again this should not
            be 0.

            If parameter {\bf reset} has an argument it provides a means of
            resetting (clearing) the output queue. The argument may have a
            different clock domain to that of the write side of the queue and hence
            this may be more convenient than directly resetting the queue where a
            change of clock domain may need to be explicitly provided.

            Note that a write to a queue variable from {\bf cpu\_api.write\_queue()} (or
            indeed from anywhere) can be detected within the FPGA by means of the
            inbuilt function {\bf waswritten()} \autorefph{sec:ifwaswritten}. When
            using this function the clock domain must be the same as that of the
            queue write, i.e. {\bf bus\_clk}.
            Parameter {\bf out} is the queue-mode output variable.
            
            On running info2c on the .json file the following C functions will have
            been defined for interface {\em name} -
            
            \begin{tabular}{| l | l | l |}
            \hline
            function                                 & when                   &                            \\ \hline \hline
            fpga\_{\em name}\_write()                & always                 & write a datum              \\ \hline
            fpga\_{\em name}\_avail\_read()          & freecount has argument & read number of spaces free \\ \hline
            fpga\_{\em name}\_avail\_thresh\_write() & ilevel "log" true      & set interrupt threshold    \\ \hline
            fpga\_{\em name}\_enable()               & ilevel has argument    & enable interrupt           \\ \hline
            fpga\_{\em name}\_disable()              & ilevel has argument    & disable interrupt          \\ \hline 
            fpga\_{\em name}\_wait()                 & ilevel has argument    & wait on interrupt          \\ \hline  
            \end{tabular}


            
            Example: \\
            If 3PL source file prog.3pl contains the following \\
            \begin{lstlisting}[gobble=12, language=threepl]
               class(cpu, cpu_api());
               queue(p, "uint:16");
               .
               .
               cpu.write_queue(p <- "P2",, true);
               .
               .
            \end{lstlisting}
            then for zynq the CPU code should contain something like \\
            \begin{lstlisting}[gobble=12, language=C]
                   #include "prog_fpga.h"
                   .
                   .
                   fpga_t *fpga = NULL;
                   int     d[...];
                   int     i = 0;
                   int     n = sizeof(d) / sizeof int;
                   .
                   .
                   fpga_open(portName, INTERRUPTS);
                   .
                   .
                   fpga_P2_avail_thresh_write(fpga, n); // set interrupt level
                   fpga_P2_wait(fpga);                  // wait for interrupt
                   for (i=0 ; i<n ; i++) {              // transmit data block
                       fpga_P2_write(fpga, d[i]);       // write item
                   fpga_P2_enable(fpga);                // enable interrupt
                   .
            \end{lstlisting}
            
            or perhaps -
            \begin{lstlisting}[gobble=12, language=threepl]
               class(cpu, cpu_api());
               queue(p, "uint:16");
               .
               .
               cpu.write_queue(p <- "P2", false);
               .
               .
            \end{lstlisting}
            then for zynq the CPU code should contain something like \\
            \begin{lstlisting}[gobble=12, language=C]
                   #include "prog_fpga.h"
                   .
                   .
                   fpga_t *fpga = NULL;
                   int     d[...];
                   int     i = 0;
                   int     n = sizeof(d) / sizeof int;
                   .
                   .
                   fpga_open(portName, INTERRUPTS);
                   .
                   .
                   for ( i=0 ; i<n ; i++) {                // transmit data block
                       while (fpga_P2_avail(fpga) == 0)    // poll until not full
                           ;
                       fpga_P2_write(fpga, d[i]);          // write item
                   }
                   .
            \end{lstlisting}

            This procedure creates no in-line target code so the current clock domain
            when this procedure is called is not relevant. The clock domain is not
            changed by this procedure,
            
            {\bf The support CPU software for the UART CPU interface is still
            under development, but application code wil be similar to the
            above.}
